\chapter{实验与分析}
\label{chap:experiment}

\section{数据集}

我们从新闻、科普、小说与技术文档四类来源中各选取 300 篇文章，
每篇 3\,000～12\,000 字。312 位读者在三种行长、两种行距下阅读，
每读完一篇回答五道理解题，并对疲劳程度打分（1～7）。

\section{整体结果}

表~\ref{tab:main} 比较了三种模型与人工评分的相关系数。
加入版面负荷后，相关系数从 $0.69$ 提高到 $0.81$。

\begin{table}[htbp]
  \centering
  \caption{不同模型与人工评分的相关系数}
  \label{tab:main}
  \begin{tabular}{@{}l c c c@{}}
    \toprule
    模型 & 新闻与科普 & 小说 & 技术文档 \\
    \midrule
    只用语言特征 & 0.71 & 0.66 & 0.70 \\
    只用版面特征 & 0.38 & 0.45 & 0.31 \\
    \textbf{二者融合（本文）} & \textbf{0.82} & \textbf{0.80} & \textbf{0.81} \\
    \bottomrule
  \end{tabular}
\end{table}

\section{行长的影响}

图~\ref{fig:fatigue} 给出了不同行长下读者自评的疲劳程度。
行长在 20～26 字时疲劳最低；超过 34 字后，疲劳随行长近似线性上升。

\begin{figure}[htbp]
  \centering
  \begin{tikzpicture}
    \begin{axis}[
      width=.8\linewidth, height=5.6cm,
      xlabel={行长（字）}, ylabel={疲劳自评（1--7）},
      xmin=12, xmax=46, ymin=2, ymax=5.6,
      axis lines=left, axis line style={mist}, tick style={mist},
      tick label style={font=\small}, label style={font=\small},
      legend style={draw=none, font=\small, at={(0.03,0.97)}, anchor=north west},
    ]
      \addplot[accent, very thick, smooth, mark=*, mark size=1.6pt] coordinates
        {(14,3.6) (18,3.0) (22,2.7) (26,2.8) (30,3.2) (34,3.7) (38,4.3) (42,4.9)};
      \addlegendentry{行距 1.8}
      \addplot[ink, very thick, smooth, dashed, mark=square*, mark size=1.4pt] coordinates
        {(14,3.8) (18,3.3) (22,3.0) (26,3.2) (30,3.7) (34,4.3) (38,4.9) (42,5.4)};
      \addlegendentry{行距 1.4}
    \end{axis}
  \end{tikzpicture}
  \caption{行长与疲劳程度}
  \label{fig:fatigue}
\end{figure}

\section{小结}

版面并非可读性的背景，而是它的一部分。对于技术文档这类信息密集的长文本，
合适的行长与行距带来的改善，几乎与把文字改写得更通俗一样多。
